\documentclass[11pt]{article}
\usepackage[margin=1in]{geometry}
\usepackage{enumitem}
% =============================================================================
% Preambles/header.tex — shared LaTeX/Beamer preamble for this project
%
% Use from a Beamer lecture by prepending:
%     \input{header}
% and compiling with TEXINPUTS=../Preambles:$TEXINPUTS (the /compile-latex
% skill does this automatically).
%
% PALETTE CONTRACT. Color names defined here MUST match the names in
% Quarto/theme-template.scss so Beamer and Quarto renderings use the same
% palette. The scripts/check-palette-sync.sh script enforces this.
%
% To customize: edit the HEX values below AND the matching $variable in
% Quarto/theme-template.scss, then run scripts/check-palette-sync.sh.
% =============================================================================

% -----------------------------------------------------------------------------
% Engine detection + encoding
%
% Our /compile-latex skill uses XeLaTeX, where inputenc is unnecessary and
% can produce encoding warnings. Only load inputenc under pdfTeX.
% -----------------------------------------------------------------------------
\usepackage{iftex}
\ifPDFTeX
  \usepackage[utf8]{inputenc}
\fi

\usepackage{amsmath, amssymb, amsthm}
\usepackage{booktabs}
\usepackage{graphicx}

% -----------------------------------------------------------------------------
% PALETTE — mirrors Quarto/theme-template.scss variable names.
% Load xcolor and define colors BEFORE anything that references them
% (notably hyperref, hypersetup, and the Beamer color assignments).
% -----------------------------------------------------------------------------
\usepackage{xcolor}

\definecolor{primary-blue}{HTML}{012169}      % headings, accents
\definecolor{primary-gold}{HTML}{B9975B}      % emphasis, borders
\definecolor{highlight-yellow}{HTML}{F2A900}  % markers, alerts
\definecolor{light-bg}{HTML}{E8EDF5}          % title / section backgrounds
\definecolor{jet}{HTML}{1A1A1A}               % body text

% Semantic colors (used in diagrams and annotations)
\definecolor{positive}{HTML}{15803D}          % good / observed / identified
\definecolor{negative}{HTML}{B91C1C}          % bad / problematic / bias
\definecolor{neutral}{HTML}{525252}           % reference / context

% Highlight accents (match .hi-* classes in the SCSS)
\definecolor{hi-slate}{HTML}{314F4F}
\definecolor{hi-green}{HTML}{2E8B57}
\definecolor{hi-red}{HTML}{C41E3A}

% Semantic aliases so skill templates and snippets can write
% \textcolor{accent}{...} without hard-coding "primary-blue".
\colorlet{accent}{primary-blue}
\colorlet{accent2}{primary-gold}

% -----------------------------------------------------------------------------
% Hyperref — loaded AFTER xcolor + palette so link colors resolve.
% -----------------------------------------------------------------------------
\usepackage{hyperref}
\hypersetup{colorlinks=true, linkcolor=primary-blue, urlcolor=primary-blue,
            citecolor=primary-blue, pdfborder={0 0 0}}

% -----------------------------------------------------------------------------
% Course code — set once per deck (after \input{header}, before \begin{document})
% via \coursecode{CS401}. Shown in the footer on every slide so a deck's course
% is identifiable without opening the title page. Empty by default.
% -----------------------------------------------------------------------------
\makeatletter
\newcommand{\coursecode}[1]{\gdef\@coursecodetext{#1}}
\gdef\@coursecodetext{}
\makeatother

% -----------------------------------------------------------------------------
% Beamer theme assignments (only apply under Beamer).
%
% We detect Beamer via \@ifclassloaded{beamer}, which is the documented,
% non-internal way. This must be wrapped in \makeatletter/\makeatother
% because \@ifclassloaded contains an @-catcode character.
% -----------------------------------------------------------------------------
\makeatletter
\@ifclassloaded{beamer}{%
  \usefonttheme{professionalfonts}
  \setbeamercolor{structure}{fg=primary-blue}
  \setbeamercolor{frametitle}{fg=primary-blue}
  \setbeamercolor{title}{fg=primary-blue}
  \setbeamercolor{subtitle}{fg=primary-gold}
  \setbeamercolor{author}{fg=primary-blue}
  \setbeamercolor{institute}{fg=neutral}
  \setbeamercolor{date}{fg=primary-gold}
  \setbeamercolor{itemize item}{fg=primary-blue}
  \setbeamercolor{itemize subitem}{fg=primary-gold}
  \setbeamercolor{alerted text}{fg=primary-gold}
  \setbeamercolor{block title}{fg=white, bg=primary-blue}
  \setbeamercolor{block body}{bg=light-bg}
  % Semantic block variants: block=definition/concept (blue), exampleblock=
  % worked example (green/positive), alertblock=key takeaway or caution (gold).
  % Keep to <=2 colored boxes per slide across ALL three variants (INV-7).
  \setbeamercolor{block title example}{fg=white, bg=positive}
  \setbeamercolor{block body example}{bg=light-bg}
  \setbeamercolor{block title alerted}{fg=white, bg=primary-gold}
  \setbeamercolor{block body alerted}{bg=light-bg}

  % Minimal footer: course code (left) + page X / N (right), no navigation clutter
  \setbeamertemplate{navigation symbols}{}
  \setbeamertemplate{footline}{%
    \color{neutral}\footnotesize\hspace{0.7em}\@coursecodetext\hfill
    \insertframenumber/\inserttotalframenumber\hspace{0.7em}\vspace{0.3em}}
}{}
\makeatother

% -----------------------------------------------------------------------------
% TikZ setup — libraries the snippet gallery uses
% -----------------------------------------------------------------------------
\usepackage{tikz}
\usetikzlibrary{arrows.meta, positioning, calc, decorations.pathreplacing,
                fit, shapes.geometric, backgrounds}

% Shared TikZ styles — snippets and hand-written diagrams can pick these up
% via `\begin{tikzpicture}[every node/.style={...}]` or by naming specific
% styles (e.g., `\node[dag-node] (X) at (0,0) {X};`).
\tikzset{
  % Node styles for causal DAGs and similar
  dag-node/.style = {
    draw=primary-blue, fill=light-bg, rounded corners=2pt,
    minimum width=1.2cm, minimum height=0.9cm, align=center, font=\small
  },
  decision-node/.style = {
    draw=primary-gold, fill=white, diamond, aspect=1.8,
    minimum width=1.8cm, minimum height=1.2cm, align=center, font=\small,
    inner sep=1pt
  },
  % Edge styles — observed vs counterfactual
  observed-edge/.style = {-{Latex[length=2.5mm]}, thick, draw=jet},
  counterfactual-edge/.style = {-{Latex[length=2.5mm]}, thick, draw=neutral, dashed},
  confound-edge/.style = {-{Latex[length=2.5mm]}, thick, draw=negative, dashed},
  % Data-point styles for plots
  observed-dot/.style = {fill=primary-blue, draw=primary-blue, circle, inner sep=1.2pt},
  counterfactual-dot/.style = {fill=white, draw=neutral, circle, inner sep=1.2pt}
}

% -----------------------------------------------------------------------------
% Convenience macros
% -----------------------------------------------------------------------------
\newcommand{\muted}[1]{\textcolor{neutral}{#1}}
\newcommand{\key}[1]{\textcolor{primary-gold}{\textbf{#1}}}
\newcommand{\good}[1]{\textcolor{positive}{#1}}
\newcommand{\bad}[1]{\textcolor{negative}{#1}}

% Standout frame macro: full-bleed dark background for section transitions.
% Beamer-only (uses \begin{frame}/\setbeamercolor/\usebeamerfont) -- do not
% call from an article-class file (e.g. Notes/<CODE>/*-notes.tex). \muted,
% \key, \good, \bad, and \coursecode above are plain xcolor/gdef and are
% safe to use from either class; verified when Notes/article-class support
% was added.
% Expand: \transitionslide{Your transition title}
\newcommand{\transitionslide}[1]{%
  \begingroup
  \setbeamercolor{background canvas}{bg=primary-blue}
  \begin{frame}[plain]
    \centering
    \vfill
    {\usebeamerfont{title}\color{white}\Large #1\par}
    \vfill
  \end{frame}
  \endgroup
}

% Section divider frame (Beamer-only), an alternative to \transitionslide with a
% darker two-line style: near-black (jet) background, a small white label line
% above a larger gold title, and a thin gold rule along the bottom edge.
% Expand: \sectiondivider{Section 2}{Pointers}
\newcommand{\sectiondivider}[2]{%
  \begingroup
  \setbeamercolor{background canvas}{bg=jet}
  \begin{frame}[plain]
    \vspace*{3.0cm}
    {\color{white}\ttfamily\large #1\par}
    \vspace{0.5cm}
    {\color{primary-gold}\LARGE #2\par}
    \begin{tikzpicture}[remember picture, overlay]
      \draw[primary-gold, line width=2.5pt]
        ($(current page.south west)+(0,0.1)$) -- ($(current page.south east)+(0,0.1)$);
    \end{tikzpicture}
  \end{frame}
  \endgroup
}

% -----------------------------------------------------------------------------
% Channel watermark — "Concepts That Click" (YouTube).
%
% Article-class files (CompetitiveExam Q/A sets, Notes, Assignments) get a
% light diagonal draft watermark on every page plus a centered footer naming
% the channel. Beamer decks get a subtle TikZ background watermark instead
% (the footline already carries the course code). Centralized here so every
% MCQ PDF is branded without per-file edits.
% -----------------------------------------------------------------------------
\makeatletter
\@ifclassloaded{article}{%
  \usepackage{draftwatermark}
  \SetWatermarkText{Concepts That Click}
  \SetWatermarkScale{1}
  \SetWatermarkFontSize{2cm}
  \SetWatermarkLightness{0.86}
  \SetWatermarkAngle{45}
  \usepackage{fancyhdr}
  \pagestyle{fancy}
  \fancyhf{}
  \fancyfoot[C]{\footnotesize\textcolor{neutral}{Concepts That Click~|~YouTube: \href{https://www.youtube.com/@concepts.thatclick}{@concepts.thatclick}}}
  \renewcommand{\headrulewidth}{0pt}
  \renewcommand{\footrulewidth}{0pt}
  \fancypagestyle{plain}{%
    \fancyhf{}%
    \fancyfoot[C]{\footnotesize\textcolor{neutral}{Concepts That Click~|~YouTube: \href{https://www.youtube.com/@concepts.thatclick}{@concepts.thatclick}}}%
    \renewcommand{\headrulewidth}{0pt}%
    \renewcommand{\footrulewidth}{0pt}%
  }
}{}
\@ifclassloaded{beamer}{%
  \setbeamertemplate{background}{%
    \begin{tikzpicture}[remember picture, overlay]
      \node[rotate=45, opacity=0.055, align=center]
        at (current page.center)
        {\Huge\textcolor{neutral}{Concepts That Click}};
    \end{tikzpicture}%
  }
}{}
\makeatother 
\coursecode{JKSSB}
\renewcommand{\thesection}{25.\arabic{section}}
\newtheorem{example}{Example}[section]
\newenvironment{definitionbox}[1]{%
  \par\noindent\textbf{Definition (#1).}\ \itshape
}{\par\vspace{0.3em}}
\title{Excel Formulas --- Detailed Lecture Notes}
\author{JKSSB Computer Awareness}
\date{\today}

\begin{document}
\maketitle

\section{What a formula is}
A \textbf{formula} is an instruction that tells Microsoft Excel how to
calculate the value of a cell. Instead of holding a fixed number, the cell
holds a rule, and Excel works out the result of that rule from the other
cells, from constants, or from functions the rule refers to. The formula is
typed into the formula bar; the result is shown in the cell. When any cell
that the formula uses changes, Excel recalculates and shows the new result
automatically.

\begin{definitionbox}{Formula}
An instruction, written in a cell or the formula bar, that calculates a value
from operands such as cell references, constants, and functions. The cell
stores the formula, while the sheet displays the result.
\end{definitionbox}

This lesson covers the four building blocks of a formula: the leading
\texttt{=}, the operators, the cell references, and the way references behave
when a formula is copied. It closes with the standard error codes and the
distinctions that examiners test most often.

\section{The \texttt{=} rule and formula syntax}
Every Excel formula must begin with the equal sign \texttt{=}. This is the
single fact most often tested on this topic. Typing \texttt{=A1+A2} into a
cell is treated as a formula and calculated; typing \texttt{A1+A2} without the
leading \texttt{=} is treated as ordinary text and shown exactly as written.

The same \texttt{=} symbol also serves as a comparison operator that tests
whether two values are equal. The position tells you which role it plays: at
the very start of an entry it begins a formula, while inside a formula it
compares two operands.

\begin{definitionbox}{The \texttt{=} rule}
A formula must start with \texttt{=}. A cell entry that does not start with
\texttt{=} is not a formula.
\end{definitionbox}

A simple formula has three parts: the leading \texttt{=}, the operands, and
the operators. In the formula \texttt{=A1*B1+100}, the operands are the
references \texttt{A1} and \texttt{B1} and the constant \texttt{100}, while
the operators are \texttt{*} (multiply) and \texttt{+} (add).

\section{Formula, function, and constant}
It helps to separate three ideas that often appear together. A
\textbf{constant} is a fixed value written directly into a formula, such as
the number \texttt{100} or a piece of text in quotation marks. A
\textbf{function} is a built-in calculation with a name, such as
\texttt{SUM} or \texttt{AVERAGE}, which is called from inside a formula. A
\textbf{formula} is the whole instruction, and it may combine constants,
references, operators, and functions in one expression.

\begin{definitionbox}{Function}
A named, built-in calculation such as \texttt{SUM(A1:A10)} or
\texttt{AVERAGE(B1:B5)} that is called from within a formula and returns a
single value.
\end{definitionbox}

The distinction matters for exams because a question may ask what a formula
contains rather than what it produces. The formula \texttt{=A1+A2} contains
two relative references and one arithmetic operator. The formula
\texttt{=SUM(A1:A10)} calls a function on a range. Both begin with \texttt{=},
and both return a value to the cell that holds them. The rest of this lesson
concentrates on the references and operators that every formula uses.

\section{Operators}
An operator tells Excel what to do with the values on either side of it.
Excel has four families of operators: arithmetic, comparison, text, and
reference.

\textbf{Arithmetic operators} produce numbers.

\begin{center}
\begin{tabular}{p{0.12\textwidth}p{0.72\textwidth}}
\toprule
\textbf{Operator} & \textbf{Meaning} \\
\midrule
\texttt{+} & Addition \\
\texttt{-} & Subtraction \\
\texttt{*} & Multiplication \\
\texttt{/} & Division \\
\texttt{\textasciicircum} & Exponent (raise to a power) \\
\texttt{\%} & Percent \\
\bottomrule
\end{tabular}
\end{center}

Two symbols are easy to get wrong. Multiplication uses the asterisk
\texttt{*}, not the letter \texttt{x} or a dot; division uses the forward
slash \texttt{/}. The caret \texttt{\textasciicircum} raises a number to a
power, so \texttt{=2\textasciicircum3} is 8. The percent sign \texttt{\%}
divides by 100, so \texttt{=50\%} is 0.5.

\textbf{Comparison operators} compare two values and return \texttt{TRUE} or
\texttt{FALSE}: \texttt{=} (equal to), \texttt{>} (greater than), \texttt{<}
(less than), \texttt{>=} (greater than or equal to), \texttt{<=} (less than or
equal to), and \texttt{<>} (not equal to). The operator \texttt{<>} is the
Excel way of writing ``not equal''; it is not written \texttt{!=}.

A comparison gives a useful result even on its own: if A1 holds 50, the
formula \texttt{=A1>100} shows \texttt{FALSE}. The operators that candidates
confuse most are \texttt{=}, which tests equality when it appears inside a
formula, and \texttt{<>}, which tests inequality. Note that the leading
\texttt{=} that starts every formula and the \texttt{=} comparison operator
are the same character used in two different roles.

\textbf{Text operator.} The ampersand \texttt{\&} joins, or concatenates, two
pieces of text into one. \texttt{="Total: "\&A1} shows the words
\texttt{Total:} followed by the contents of cell A1.

\textbf{Reference operators} build references to groups of cells. The colon
\texttt{:} is the range operator, so \texttt{A1:A10} means every cell from A1
down to A10. The comma \texttt{,} is the union operator, which joins two
ranges into one list. A single space between two references is the
intersection operator, which returns the cells common to both ranges.

The four families can be summarised in one place. Arithmetic operators
(\texttt{+ - * / \textasciicircum \%}) produce numbers. Comparison operators
(\texttt{= > < >= <= <>}) produce \texttt{TRUE} or \texttt{FALSE}. The text
operator \texttt{\&} produces a joined piece of text. The reference operators
\texttt{:}, \texttt{,}, and the space build ranges and lists of cells.
Knowing which family an operator belongs to tells you what kind of result to
expect.

\section{Order of evaluation (precedence)}
When a formula has several operators, Excel does not simply work from left to
right. It follows a fixed order of precedence. Parentheses are evaluated
first. Then the percent sign, then the exponent, then multiplication and
division, then addition and subtraction. After the arithmetic comes the
concatenation operator \texttt{\&}, and finally the comparison operators.

Because multiplication and division are evaluated before addition and
subtraction, \texttt{=2+3*4} is 14, not 20. The multiplication \texttt{3*4} is
done first to give 12, and then 2 is added. In the same way
\texttt{=2+3*4} differs from \texttt{=(2+3)*4}, which is 20. Parentheses are
the reliable way to control the order of evaluation.

\begin{example}
In \texttt{=10+2*5} the multiplication is done first, so the result is 20. If
the intention was to add first, the formula must be written
\texttt{=(10+2)*5}, which gives 60. The exponent binds even tighter than
multiplication, so \texttt{=2*3\textasciicircum2} is 18, because
\texttt{3\textasciicircum2} is evaluated before the multiplication by 2.
\end{example}

\section{Cell references}
A \textbf{cell reference} names a cell by its column letter and row number. In
the default \texttt{A1} reference style, \texttt{A1} is column A, row 1, and
\texttt{B5} is column B, row 5. A reference lets a formula use the current
value of another cell, so the formula stays correct when that value changes.

The crucial property of a reference is not what it points at now but how it
behaves when the formula is \emph{copied}. There are three kinds of
reference: relative, absolute, and mixed.

\section{Relative references}
A reference such as \texttt{A1} is \textbf{relative}. Excel stores it, in
effect, as an offset from the cell that contains the formula. When the formula
is copied to another cell, the reference shifts by the same distance, so it
keeps pointing at the cell that lies in the same relative position.

\begin{example}
If \texttt{=A1} is entered in cell B1 and then copied to B2, the formula
becomes \texttt{=A2}. The reference moved down one row because the copied
formula moved down one row. Copy the same formula from B1 to C1, and it
becomes \texttt{=B1}: the reference moved one column to the right because the
formula did.
\end{example}

Relative references are the default and are what you want for a column of
repeated calculations, such as adding each row's two numbers.

\section{Absolute references}
A reference such as \texttt{\$A\$1} is \textbf{absolute}. The dollar sign
locks the part of the reference that follows it: the first \texttt{\$} locks
the column, and the second \texttt{\$} locks the row. An absolute reference
does not change when the formula is copied; it always points at the same cell.

\begin{example}
If \texttt{=\$A\$1} is entered in B1 and copied to B5, it remains
\texttt{=\$A\$1}. Neither the column nor the row moved. A fixed value such as
a tax rate or a constant kept in one cell should be referenced absolutely, so
that copying the formula does not disturb it.
\end{example}

\section{Mixed references}
A \textbf{mixed} reference locks only one part. In \texttt{\$A1} the column is
locked and the row is relative; copying the formula to another row changes the
row number but never the column. In \texttt{A\$1} the row is locked and the
column is relative; copying the formula to another column changes the column
letter but never the row.

A useful way to remember this is that the \texttt{\$} locks whatever
immediately follows it. A \texttt{\$} before a letter locks the column, and a
\texttt{\$} before a number locks the row.

\section{Copying and filling formulas}
The behaviour of a reference under copying is the heart of the topic. When a
formula is copied, only the unlocked parts of each reference shift; the locked
parts stay put. The table below compares the four reference types when a
formula is copied one column right and one row down, from cell B1 to cell C2.

\begin{center}
\begin{tabular}{p{0.20\textwidth}p{0.30\textwidth}p{0.34\textwidth}}
\toprule
\textbf{Reference} & \textbf{What is locked} & \textbf{After copying B1 to C2} \\
\midrule
\texttt{A1} & nothing & \texttt{B2} \\
\texttt{\$A\$1} & column and row & \texttt{\$A\$1} \\
\texttt{\$A1} & column & \texttt{\$A2} \\
\texttt{A\$1} & row & \texttt{B\$1} \\
\bottomrule
\end{tabular}
\end{center}

Reading the table: the relative reference \texttt{A1} moved on both axes; the
absolute reference \texttt{\$A\$1} moved on neither; \texttt{\$A1} kept its
column but followed the row; and \texttt{A\$1} kept its row but followed the
column.

A second, very common direction is copying a formula straight down a column.
If \texttt{=A1*B1} sits in C1 and is copied down to C2, the formula becomes
\texttt{=A2*B2}: both relative references followed the copy down one row. If
the value in B1 must stay fixed, the formula should instead read
\texttt{=A1*\$B\$1}, so that only A1 follows the copy while \texttt{\$B\$1}
stays anchored.

The same four references behave as follows when copied one column to the
right, and when copied one row down.

\begin{center}
\begin{tabular}{p{0.20\textwidth}p{0.32\textwidth}p{0.32\textwidth}}
\toprule
\textbf{Reference} & \textbf{Copied right (B1$\rightarrow$C1)} & \textbf{Copied down (B1$\rightarrow$B2)} \\
\midrule
\texttt{A1} & \texttt{B1} & \texttt{A2} \\
\texttt{\$A\$1} & \texttt{\$A\$1} & \texttt{\$A\$1} \\
\texttt{\$A1} & \texttt{\$A1} & \texttt{\$A2} \\
\texttt{A\$1} & \texttt{B\$1} & \texttt{A\$1} \\
\bottomrule
\end{tabular}
\end{center}

Reading either column, the rule is the same: a reference moves only along the
axis that is not locked. The relative reference moves on both axes, the
absolute reference on neither, and each mixed reference on exactly one.

\section{Worked examples}

\begin{example}
A price list has prices in column A and a fixed tax rate stored in cell
\texttt{\$B\$1}. The formula in cell C2 is \texttt{=A2*\$B\$1}. When it is
copied down, C3 becomes \texttt{=A3*\$B\$1} and C4 becomes \texttt{=A4*\$B\$1}:
the price reference moves down with each row, while the rate reference stays
locked on \texttt{\$B\$1}. If the formula had been written \texttt{=A2*B1},
copying it down would have moved the rate to B2 and then B3, giving wrong
results.
\end{example}

\begin{example}
In cell B1 the formula is \texttt{=\$A1}. It is copied to C2. The column lock
holds A, and the relative row follows the copy down one row, so C2 contains
\texttt{=\$A2}. If instead the original formula were \texttt{=A\$1}, copying
from B1 to C2 would keep the row 1 and move the column one step from A to B,
giving \texttt{=B\$1}.
\end{example}

\begin{example}
A sheet lists marks in column B. To flag each student as passing, the formula
\texttt{=B2>=40} is entered in C2 and copied down. It returns \texttt{TRUE}
for a mark of 40 or more and \texttt{FALSE} otherwise. Because B2 is relative,
the reference follows the copy down the column, so C3 tests B3 and C4 tests
B4. The comparison operator produces the \texttt{TRUE}/\texttt{FALSE} result,
not a number.
\end{example}

\begin{example}
A score sheet holds marks in B2:B11 and the class total in cell \texttt{\$B\$12}.
To show each student's share as a percentage, the formula in C2 is
\texttt{=B2/\$B\$12*100}. The relative \texttt{B2} follows the copy down the
column, while \texttt{\$B\$12} stays anchored on the total. The operators here
are division and multiplication, evaluated from left to right at the same
precedence, and the result is a plain number. If the total were written as a
relative reference, copying the formula down would corrupt the denominator.
\end{example}

\section{Common formula errors}
When a formula cannot be evaluated, Excel shows an error value in the cell.
The codes tested in exams are:

\begin{center}
\begin{tabular}{p{0.18\textwidth}p{0.68\textwidth}}
\toprule
\textbf{Error} & \textbf{Meaning} \\
\midrule
\texttt{\#DIV/0!} & Division by zero, or by an empty cell. \\
\texttt{\#VALUE!} & The wrong type of operand, such as text where a number is needed. \\
\texttt{\#REF!} & A referenced cell was deleted, so the reference is no longer valid. \\
\texttt{\#NAME?} & An unknown function, or a misspelled name or text entry. \\
\texttt{\#N/A} & The value is not available, as when a lookup finds no match. \\
\texttt{\#NULL!} & The two ranges given do not intersect. \\
\bottomrule
\end{tabular}
\end{center}

Remember the two that examiners use most: \texttt{\#DIV/0!} for division by
zero and \texttt{\#REF!} for a reference that has been broken by deleting a
cell. A related warning is the \textbf{circular reference}, which appears when
a formula refers, directly or through a chain of other formulas, back to its
own cell, so Excel cannot work out a value. Deleting the referenced row or
column is the usual cause of \texttt{\#REF!}, while supplying the wrong type
of operand, such as text where a number is expected, produces
\texttt{\#VALUE!}.

\section{Exam retrieval and common confusions}
Six distinctions prevent most errors on this topic.
\begin{enumerate}
  \item A formula must start with \textbf{\texttt{=}} (TRAP-30, PYQ-38). The
    options \texttt{+}, \texttt{\%}, and \texttt{@} are distractors.
  \item \textbf{\texttt{\$}} before a letter locks the column; \textbf{\texttt{\$}}
    before a number locks the row. Do not swap these.
  \item \textbf{Relative} references change on copy; \textbf{absolute}
    references do not. A \textbf{mixed} reference changes on one axis only.
  \item \texttt{\#REF!} means a broken reference; \texttt{\#DIV/0!} means
    division by zero. Do not swap the error codes.
  \item \texttt{<>} is the Excel operator for ``not equal''; multiplication is
    \texttt{*} and division is \texttt{/}.
  \item A \textbf{function} such as \texttt{SUM} is called from inside a
    formula, and a \textbf{constant} is a fixed value written directly; both
    are operands, not operators.
\end{enumerate}

\begin{example}
An item asks what \texttt{=\$A\$1} becomes when copied from B1 to B5. Because
both the column and the row are locked by the two dollar signs, the reference
does not move: the answer is \texttt{=\$A\$1}. If the question instead gives
\texttt{=\$A1}, the column is locked but the row is relative, so the answer is
\texttt{=\$A5}. The dollar signs, not the row numbers, decide the answer.
\end{example}

\subsection*{Exercises}
\begin{enumerate}[label=\arabic*.]
  \item Which character must every Excel formula begin with, and what happens
    if it is omitted?
  \item List the arithmetic operators for addition, subtraction,
    multiplication, division, and exponent.
  \item Explain the difference between a relative reference \texttt{A1} and an
    absolute reference \texttt{\$A\$1}.
  \item In the mixed reference \texttt{\$A1}, which part is locked and which
    part changes when the formula is copied?
  \item The formula \texttt{=A1} is in cell B1. What does B2 contain after the
    formula is copied to it? What would B2 contain if the formula were
    \texttt{=\$A\$1}?
  \item State the meaning of the errors \texttt{\#DIV/0!} and
    \texttt{\#REF!}.
  \item What is the result of the formula \texttt{=2+3*4}, and why?
  \item Name the Excel operator that joins two pieces of text, and the operator
    that means ``not equal to''.
\end{enumerate}

\subsection*{Solutions}
\begingroup\small
\begin{enumerate}[label=\arabic*.]
  \item A formula must begin with \texttt{=}. Without it, the entry is stored
    and shown as ordinary text rather than calculated.
  \item Addition \texttt{+}, subtraction \texttt{-}, multiplication
    \texttt{*}, division \texttt{/}, and exponent \texttt{\textasciicircum}.
  \item \texttt{A1} is relative: on copying, both its column and its row shift
    by the same offset as the formula. \texttt{\$A\$1} is absolute: the dollar
    signs lock the column and the row, so it does not change on copying.
  \item In \texttt{\$A1} the column A is locked; the row is relative, so the
    row number changes when the formula is copied to a different row.
  \item Copying \texttt{=A1} from B1 to B2 gives \texttt{=A2}, because the
    relative reference follows the copy down one row. Copying \texttt{=\$A\$1}
    to B2 leaves it as \texttt{=\$A\$1}, because both parts are locked.
  \item \texttt{\#DIV/0!} means division by zero or by an empty cell;
    \texttt{\#REF!} means a reference is invalid because a referenced cell was
    deleted.
  \item The result is 14. Multiplication is evaluated before addition, so
    \texttt{3*4} is 12, and \texttt{2+12} is 14.
  \item The ampersand \texttt{\&} joins (concatenates) text, and \texttt{<>}
    means ``not equal to''.
\end{enumerate}
\endgroup
\end{document}
